\chapter{最后要看的作业}
\section{第一章作业}
\subsection*{1.1-1.3}

\begin{homework}[1.1]
写出下列随机试验的样本空间：

\begin{enumerate}
  \item 掷三颗骰子，观察所得点数之和；
  \item 抽检产品，直到抽到 $10$ 件正品为止，记录抽检产品的数量；
  \item 从编号分别为 $1,2,3,4,5$ 的 $5$ 件产品中任取两件，观察取出哪两件；
  \item 观察某地一天内的最高气温和最低气温（假设最低气温不低于 $T_1$，最高气温不高于 $T_2$）；
  \item 将长为 $l$ 的线段分成三段，观察各段长度。
\end{enumerate}
\end{homework}

\begin{homework}[1.2]
设 $A,B,C$ 为事件，用 $A,B,C$ 的运算关系表示下列事件：

\begin{enumerate}
  \item $A$ 与 $B$ 出现，但 $C$ 不出现；
  \item $A$ 出现，且 $B$ 与 $C$ 至少一个出现；
  \item $A,B,C$ 中至少一个出现；
  \item $A,B,C$ 中恰有一个出现；
  \item $A,B,C$ 中至少两个出现；
  \item $A,B,C$ 中恰有两个出现；
  \item $A,B,C$ 中至少多个出现；
  \item $A,B,C$ 中至少多个两个出现。
\end{enumerate}
\end{homework}

\begin{homework}[1.12]
甲、乙两货轮驶向一个不能使它们同时使用的码头。假设两货轮在一昼夜内各时刻到达码头都是等可能的，且甲货轮使用码头的时间是 $1$ 小时，乙货轮使用码头的时间是 $2$ 小时，求两货轮到达码头后能直接使用码头的概率。
\end{homework}

\begin{homework}[1.13]
将长为 $l$ 的线段折成三段，求下列事件的概率：

\begin{enumerate}
  \item 三段可构成一个三角形；
  \item 最长一段不超过 $2l/3$。
\end{enumerate}
\end{homework}

\begin{homework}[1.14]
在区间 $(0,1)$ 内任取两数，求下列事件的概率：

\begin{enumerate}
  \item 两数之差大于 $0.2$；
  \item 两数之和小于 $1.2$；
  \item 前两个事件同时出现。
\end{enumerate}
\end{homework}

\begin{homework}[1.15]
设 $A,B$ 是事件，$P(A)=x,\ P(B)=y,\ P(AB)=z$。求：

\begin{enumerate}
  \item $P(\overline{A \cup B})$；
  \item $P(\overline{AB})$；
  \item $P(\overline{A} \cup B)$；
  \item $P(\overline{A}\ \overline{B})$。
\end{enumerate}
\end{homework}

\begin{homework}[1.16]
设 $A,B,C$ 是事件，$P(A)=P(B)=P(C)=0.25,\ P(AB)=P(BC)=0,\ P(AC)=0.125$，求 $A,B,C$ 中至少一个出现的概率。
\end{homework}

\begin{homework}[1.17]
一辆载有 $25$ 名乘客的巴士从机场开出，沿途经 $9$ 个站点。假设每位乘客在任一站点下车是等可能的，且只有在有乘客下车时巴士才停车。求下列事件的概率：

\begin{enumerate}
  \item 巴士在第 $i$ 站停车，$i=1,2,\dots,9$；
  \item 巴士在第 $i$ 站和第 $j$ 站均停车，$1 \le i < j \le 9$；
  \item 巴士在第 $i$ 站和第 $j$ 站至少停 $1$ 次，$1 \le i < j \le 9$；
  \item 在第 $i$ 站有 $3$ 人下车，$i=1,2,\dots,9$。
\end{enumerate}
\end{homework}
\subsection*{1.4}
\begin{homework}[1.18]
设 $A,B$ 是事件，$P(A)=0.3,\ P(B)=0.4,\ P(AB)=0.5$，求 $P(B \mid A \cup \overline{B})$。
\end{homework}

\begin{homework}[1.19]
设 $A,B$ 是事件，$P(A)=1/4,\ P(B \mid A)=1/3,\ P(A \mid B)=1/2$，求 $P(A \cup B)$。
\end{homework}

\begin{homework}[1.20]
设 $A,B$ 是事件，$P(A)=a,\ P(B)=b>0$，试证：
\[
P(A \mid B) \ge \frac{a+b-1}{b}.
\]
\end{homework}

\begin{homework}[1.21]
以往数据显示，某家庭患感冒规律如下：孩子感冒的概率为 $0.6$；孩子感冒的条件下，母亲感冒的概率为 $0.5$；孩子与母亲感冒的条件下，父亲感冒的概率为 $0.4$。求孩子和母亲感冒，但父亲未感冒的概率。
\end{homework}

\begin{homework}[1.22]
$10$ 件产品中有 $2$ 件次品、$8$ 件正品。作不返回抽取，每次任取 $1$ 件，取 $2$ 次，求下列事件的概率：

\begin{enumerate}
  \item 取 $2$ 件正品；
  \item 取 $2$ 件次品；
  \item 第二次取到次品。
\end{enumerate}
\end{homework}

\begin{homework}[1.24]
某仪器有三个指示灯，其工作相互独立，烧坏第一、第二、第三个指示灯的概率分别为 $0.1, 0.2$ 和 $0.3$。当烧坏 $1$ 个灯时，仪器出现故障的概率为 $0.25$；当烧坏 $2$ 个灯时，仪器出现故障的概率为 $0.6$；当烧坏 $3$ 个灯时，仪器出现故障的概率为 $0.9$。求仪器出现故障的概率。
\end{homework}

\begin{homework}[1.25]
某地区男女两性人口比例为 $51:49$，男性中的 $5\%$ 是色盲患者，女性中的 $2.5\%$ 是色盲患者。今从人群中随机地抽取一人，恰好是色盲患者，求此人为男性的概率。
\end{homework}

\begin{homework}[1.27]
设一架坠毁的飞机掉在 $3$ 个可能区域中的任何一个是等可能的；如果飞机坠毁在区域 $A_i\ (i=1,2,3)$，经快速搜索后能发现其残骸的概率为 $\alpha_i\ (i=1,2,3)$。现已快速搜寻了区域 $A_1$，但未发现飞机残骸。求飞机坠落在 $3$ 个区域内的条件概率。
\end{homework}
\begin{dxtips}[注意了谁在前面谁在后面最后要演算一下]
    \begin{solution}
设事件 $B$ 表示“已快速搜寻区域 $A_1$ 但未发现残骸”。题设知
\[
P(A_i)=\frac13,\quad i=1,2,3,
\]
且
\[
P(B\mid A_1)=1-\alpha_1,\qquad P(B\mid A_2)=1,\qquad P(B\mid A_3)=1.
\]
由全概率公式，
\[
P(B)=\sum_{i=1}^3 P(B\mid A_i)P(A_i)
=\frac13(1-\alpha_1)+\frac13+\frac13
=1-\frac{\alpha_1}{3}.
\]
由贝叶斯公式，
\[
P(A_1\mid B)=\frac{P(B\mid A_1)P(A_1)}{P(B)}
=\frac{(1-\alpha_1)\cdot \frac13}{1-\frac{\alpha_1}{3}}
=\frac{1-\alpha_1}{3-\alpha_1},
\]
\[
P(A_2\mid B)=\frac{P(B\mid A_2)P(A_2)}{P(B)}
=\frac{1\cdot \frac13}{1-\frac{\alpha_1}{3}}
=\frac{1}{3-\alpha_1},\qquad
P(A_3\mid B)=\frac{1}{3-\alpha_1}.
\]
\end{solution}
\end{dxtips}
\begin{homework}[1.28]
设某种昆虫产 $n$ 个卵的概率为 $e^{-\lambda}\lambda^n/n!,\ n=0,1,2,\dots$, $\lambda>0$ 为常数，每个卵能孵化成幼虫的概率均为 $p(0<p<1)$，且各卵能否孵化成幼虫相互独立。

\begin{enumerate}
  \item 求每只昆虫有 $k$ 个后代（幼虫）的概率，$k=0,1,2,\dots$；
  \item 若已知某只昆虫有 $k$ 个后代，求它曾产了 $m$ 个卵的概率，$m=k,k+1,\dots$。
\end{enumerate}
\end{homework}
\section{第二章}
\subsection*{2.1}
\begin{dxtips}
    把因式分解好方便看什么消了
\end{dxtips}
\begin{homework}[2.3]
$15$ 件同型号零件中恰有 $2$ 件次品。作不放回抽取，取 $3$ 次，每次取 $1$ 件，以 $X$ 表示 $3$ 次中取出次品的件数，求 $X$ 的概率分布。
\end{homework}
\begin{dxtips}
    确定需要完成多少次，不确定总次数问题
\end{dxtips}
\begin{dxtips}
    直到成功一次问题就是n-1次不成功+最后一次一定成功
\end{dxtips}
\begin{homework}[2.4]
进行重复独立试验，每次试验成功的概率为 $p \ (0 < p < 1)$，失败的概率为 $q = 1 - p$。  
(1) 试验进行到出现 $1$ 次成功为止，$X$ 表示所需的试验次数，求 $X$ 的概率分布；  
(2) 试验进行到出现 $r$ 次成功为止，$Y$ 表示所需的试验次数，求 $Y$ 的概率分布。
\end{homework}

\begin{homework}[2.5]
甲、乙轮流掷投篮，直至某人投中为止。甲的命中率为 $0.4$，乙的命中率为 $0.6$，甲先投。分别求甲、乙投篮次数的概率分布。
\end{homework}

\begin{homework}[2.9]
甲、乙投篮的命中率分别为 $0.6$ 与 $0.7$，各投 $3$ 次。求下列事件的概率：  
(1) 甲、乙投中次数相等；  
(2) 甲投中次数多于乙投中次数。
\end{homework}
\begin{dxtips}
泊松分布
\end{dxtips}
\begin{homework}[2.12]
为保证公司设备的有效运行，需要配备一定数量的设备维修人员。假设公司有同类设备 $180$ 台，且各台设备的工作相互独立，任一时刻设备发生故障的概率都是 $0.01$。假设一台设备的故障需由一人进行修理，问至少应配备多少名维修人员，才能保证设备发生故障后能够得到及时修理的概率不小于 $0.99$？
\end{homework}

\begin{homework}[2.13]
某小学校有 $730$ 名学生，假设学生的生日等可能地出现在一年的每一天，求该校恰有 $4$ 名学生生日在同一天的概率。
\end{homework}
\subsection*{2.2}

\begin{homework}[2.15]
设连续型随机变量 $X$ 有概率密度函数
\[
f(x) = a \exp(-|x|), \quad -\infty < x < \infty,
\]
求待定常数 $a$ 及 $P(0 \leq X \leq 1)$。
\end{homework}

\begin{homework}[2.17]
顾客到达银行后一般需排队等待服务，设排队等待的时间 $X$（以分钟计）有概率密度函数
\[
f(x) =
\begin{cases}
0.2 \exp(-0.2x), & x > 0, \\[6pt]
0, & x \leq 0,
\end{cases}
\]
若某顾客每次到达银行后最多排队等待服务 $10$ 分钟，否则放弃服务而离开。他一个月到银行 $5$ 次，以 $Y$ 表示其一个月内放弃服务而离开的次数，求 $Y$ 的概率分布。
\end{homework}

\begin{homework}[2.19]
设随机变量 $X \sim N(3, 2^2)$。  
(1) 求 $P(-2 < X \leq 10)$ 和 $P(|X| > 2)$；  
(2) 确定常数 $c$，使 $P(X > c) = 0.05$。
\end{homework}
\begin{dxtips}
    次数很大但是概率不够小不能用柏松得用正态。
\end{dxtips}
\begin{homework}[2.21]
设某种元件的寿命 $X$（以小时计）服从正态分布 $N(160,\sigma^2)$。若要求
\[
P(120 < X < 200) \geq 0.80,
\]
求参数 $\sigma$ 的最大值。
\end{homework}

\begin{homework}[2.22]
某车间有 $200$ 台同型号设备，各设备工作与否相互独立。每台设备有 $60\%$ 的时间处于工作状态，工作时消耗的电能为 $W$。问该车间至少需要多少电能，才能以 $99.9\%$ 以上的概率，保证车间不因供电不足而影响生产？
\end{homework}

\begin{homework}[2.23]
设连续型随机变量 $X$ 有概率密度函数
\[
f(x) =
\begin{cases}
a \cos x, & -0.5\pi < x < 0.5\pi, \\[6pt]
0, & \text{其他},
\end{cases}
\]
其中 $a$ 是待定常数。  
(1) 求 $a$ 及 $X$ 的分布函数 $F(x)$；  
(2) 画出 $f(x)$ 和 $F(x)$ 的图形。
\end{homework}

\begin{homework}[2.24]
设 $F(x) = a \exp(-x), \; x > 0$；当 $x < 0$ 时，$F(x) = 0$。  
问 $a$ 取何值，可使 $F(x)$ 是一个分布函数？
\end{homework}
\begin{dxtips}
    \begin{itemize}
  \item $\dfrac{d}{dx}\arcsin x=\dfrac{1}{\sqrt{1-x^2}},\quad |x|<1.$
  \item $\dfrac{d}{dx}\arccos x=-\dfrac{1}{\sqrt{1-x^2}},\quad |x|<1.$
  \item $\dfrac{d}{dx}\arctan x=\dfrac{1}{1+x^2},\quad x\in\mathbb R.$
\end{itemize}
\end{dxtips}
\begin{homework}[2.25]
设连续型随机变量 $X$ 的分布函数
\[
F(x) =
\begin{cases}
0, & x < -1, \\[6pt]
a + b \arcsin x, & -1 \leq x \leq 1, \\[6pt]
1, & x > 1,
\end{cases}
\]
其中 $a$ 和 $b$ 是待定常数。求 $a$ 和 $b$。
\end{homework}
\section{第五次作业}
26--30

\begin{homework}[26]
袋中有同型号小球 $12$ 只，其中 $10$ 只红色，$2$ 只黑色。从袋中取球两次，每次任取一球。令
\[
X_i =
\begin{cases}
1, & 第 i 次取到红球, \\
0, & 第 i 次取到黑球,
\end{cases}
\quad i=1,2.
\]
试分别在放回抽取和不放回抽取情况下，求 $(X_1, X_2)$ 的概率分布和边缘概率分布。
\end{homework}

\begin{homework}[27]
设二维连续型随机变量 $(X,Y)$ 有概率密度函数
\[
f(x,y) =
\begin{cases}
a(6 - x - y), & 0 < x < 2,\; 2 < y < 4, \\[6pt]
0, & \text{其他},
\end{cases}
\]
其中 $a$ 是待定常数。求 $a$ 及 $P(X < 1, Y < 3)$。
\end{homework}

\begin{homework}[28]
设二维连续型随机变量 $(X,Y)$ 在以原点为中心、$r$ 为半径的圆内服从均匀分布，
求 $(X,Y)$ 的联合概率密度函数和边缘概率密度函数。
\end{homework}

\begin{homework}[29]
设二维连续型随机变量 $(X,Y)$ 有概率密度函数
\[
f(x,y) =
\begin{cases}
a \exp(-y), & 0 < x < y, \\[6pt]
0, & \text{其他},
\end{cases}
\]
其中 $a$ 是待定常数。求 $a$ 及边缘概率密度函数。
\end{homework}

\begin{homework}[30]
设二维连续型随机变量 $(X,Y)$ 有概率密度函数
\[
f(x,y) =
\begin{cases}
a y (2 - x), & 0 \le y \le x \le 1, \\[6pt]
0, & \text{其他},
\end{cases}
\]
其中 $a$ 是待定常数。求 $a$ 及边缘概率密度函数。
\end{homework}
\subsection{2.3}
\begin{homework}[2.31]
设某机床加工螺栓的长度 $X$（以厘米计）服从正态分布 $N(10.05, 0.06^2)$。若螺栓长度在区间 $[10.05-0.12,\; 10.05+0.12]$ 内为合格品，求该机床加工螺栓的合格品率。
\end{homework}

\begin{homework}[2.34]
设连续型随机变量 $X \sim N(0,1)$，求 $Y = |X|$ 的概率密度函数。
\end{homework}
\begin{dxtips}
    Y=sinx求的是Y的概率密度，已知的是x的概率密度能求到是x到概率分布P（x小于a），知道了P（x小于a）
    我们就可以知道P（Y小于b）
    只要把Y=sinX带入
    然后把x对应范围求出来。知道了x的范围就能积分求出P（Y<b)的具体值或者函数，然后求导就是概率密度了
\end{dxtips}
\begin{homework}[2.35]
设连续型随机变量 $X$ 有概率密度函数
\[
f_X(x) =
\begin{cases}
\dfrac{2x}{\pi^2}, & 0 < x < \pi, \\[6pt]
0, & \text{其他},
\end{cases}
\]
求 $Y = \sin X$ 的概率密度函数。
\end{homework}
\begin{dxtips}
    这里x的取值就是0到acrsinb，π-arcsinb，到π，P（Y小于b）就能算出来然后求导。
\end{dxtips}
\begin{homework}[2.37]
以 $X$ 表示某医院一天出生婴儿的个数，$Y$ 表示其中男婴的个数。设 $(X,Y)$ 有概率分布
\[
P(X = n, Y = m) = \frac{(7.14)^m (6.86)^{n-m} e^{-14}}{m!(n-m)!}, 
\quad m = 0,1,2,\ldots,n; \; n = 0,1,2,\ldots
\]
(1) 求条件概率分布；  
(2) 具体写出 $\{ P(Y = m \mid X = 20) \}$。
\end{homework}

\begin{homework}[2.39]
设二维连续型随机变量 $(X,Y)$ 有概率密度函数
\[
f(x,y) =
\begin{cases}
x^2 + \dfrac{xy}{3}, & 0 \leq x \leq 1, \; 0 \leq y \leq 2, \\[6pt]
0, & \text{其他},
\end{cases}
\]
求条件概率密度函数 $f_{X|Y}(x|y)$ 和 $f_{Y|X}(y|x)$。
\end{homework}
\subsection{2.4}
\begin{homework}[2.42]
设连续型随机变量 $X$ 有概率密度函数
\[
f_X(x) =
\begin{cases}
\dfrac{2x}{\pi^2}, & 0 < x < \pi,\\[6pt]
0, & \text{其他},
\end{cases}
\]
求 $Y = \sin X$ 的概率密度函数。
\end{homework}

\begin{homework}[2.45]
设某物体的华氏温度 $T \sim N(98.6, 2)$，求其摄氏温度
\[
Y = \tfrac{5}{9}(T - 32)
\]
的概率密度函数。
\end{homework}

\begin{homework}[2.50]
设二维连续型随机变量 $(X,Y)$ 有概率密度函数
\[
f(x,y) = \frac{1}{2\pi ab} \exp \!\left[ -\tfrac{1}{2} \!\left( \frac{x^2}{a^2} + \frac{y^2}{b^2} \right) \right],
\]
求 $(X,Y)$ 取值于椭圆
\[
\frac{x^2}{a^2} + \frac{y^2}{b^2} = r^2
\]
内的概率。
\end{homework}

\begin{homework}[2.53]
设二维随机变量 $(X,Y)$ 有习题 2.39 中的概率密度函数，求
\[
P(X+Y > 1),\qquad P(Y < X),\qquad P(Y < 0.5 \mid X < 0.5).
\]
\end{homework}

\begin{homework}[2.56]
设随机变量 $X$ 与 $Y$ 相互独立，分别服从参数为 $\lambda_1,\lambda_2$ 的泊松分布，证明 $X+Y$ 也服从泊松分布，且
\[
P(X = k \mid X+Y = N)
 = b\!\left(k;N,\frac{\lambda_1}{\lambda_1+\lambda_2}\right),
\quad k=0,1,\ldots,N.
\]
\end{homework}

\begin{homework}[2.52]
设随机变量 $X_1$ 与 $X_2$ 相互独立，分别有概率密度函数
\[
f_{X_i}(x_i)=
\begin{cases}
\lambda_i e^{-\lambda_i x_i}, & x_i > 0,\\[4pt]
0, & x_i \le 0,
\end{cases}
\quad i=1,2,
\]
其中 $\lambda_1,\lambda_2>0$。令
\[
Y =
\begin{cases}
1, & X_1 \le X_2,\\
0, & X_1 > X_2,
\end{cases}
\]
求 $Y$ 的概率分布。
\end{homework}

\begin{homework}[2.57]
设随机变量 $X_1$ 与 $X_2$ 相互独立，有概率分布
\[
P(X_i = k) = p(1-p)^k,\qquad k = 0,1,2,\ldots,\; 0<p<1,\; i=1,2.
\]
(1) 证明 $P(X_1=k \mid X_1+X_2=n)=\frac{1}{n+1}$；\\
(2) 求 $Y = \max(X_1,X_2)$ 的概率分布；\\
(3) 求 $(X_1,Y)$ 的概率分布。
\end{homework}

\begin{homework}[2.60]
设随机变量 $X_1,X_2,X_3$ 相互独立，概率密度函数均为
\[
f(x)=
\begin{cases}
\dfrac{x}{\sigma^2}\exp\!\left(-\dfrac{x^2}{2\sigma^2}\right), & x>0,\\
0, & x\le0,
\end{cases}
\]
求：  
(1) $Y=\max(X_1,X_2,X_3)$ 的概率密度函数；\\
(2) $Z=\min(X_1,X_2,X_3)$ 的概率密度函数；\\
(3) $P(Y\ge\sigma)$。
\end{homework}
\begin{dxtips}
    \[
\begin{pmatrix}dx\\dy\end{pmatrix}
=
\begin{pmatrix}
\dfrac{\partial x}{\partial u} & \dfrac{\partial x}{\partial v}\\[6pt]
\dfrac{\partial y}{\partial u} & \dfrac{\partial y}{\partial v}
\end{pmatrix}
\begin{pmatrix}du\\dv\end{pmatrix}.
\]

展开即
\[
dx=x_u\,du+x_v\,dv,\qquad
dy=y_u\,du+y_v\,dv.
\]
看概率守恒：

\begin{itemize}
  \item 在 $(x,y)$ 坐标中，小块概率近似为
  $f_{X,Y}(x,y)\,dx\,dy$；
  \item 在 $(u,v)$ 坐标中，同一小块概率也等于
  $f_{U,V}(u,v)\,du\,dv$。
\end{itemize}

因此
\[
f_{X,Y}(x,y)\,dx\,dy
=
f_{U,V}(u,v)\,du\,dv.
\]

将
$dx\,dy=|J|\,du\,dv$
代入，得到
\[
f_{X,Y}(x,y)\,|J|\,du\,dv
=
f_{U,V}(u,v)\,du\,dv.
\]
\end{dxtips}
\begin{homework}[2.61]
设随机变量 $X$ 与 $Y$ 相互独立，概率密度函数均为
\[
f(x)=
\begin{cases}
e^{-x}, & x>0,\\[4pt]
0, & x\le0,
\end{cases}
\]
(1) 证明 $X+Y$ 与 $X/Y$ 相互独立；\\
(2) 证明 $X+Y$ 与 $X/(X+Y)$ 也相互独立。
\end{homework}
\begin{dxtips}
    记住乘的是雅各比行列式的绝对值，而且在变量替换后相应的取值范围也会变比如第二小问的v取值范围就变成了0到1
\end{dxtips}
\begin{homework}[2.62]
设二维连续型随机变量 $(X,Y)$ 服从非退化二维正态分布
\[
N(0,\sigma_1^2;\;0,\sigma_2^2;\;0),
\]
求 $X-Y$ 的概率密度函数。
\end{homework}

\begin{homework}[2.63]
设连续型随机变量 $X$ 与 $Y$ 相互独立，求下列两小题中 $X-Y$ 与 $XY$ 的概率密度函数：  
(1) $X,Y$ 均服从均匀分布 $U(-a,a)$；\\
(2) $X,Y$ 均服从标准正态分布。
\end{homework}
\section{第三章}
\subsection*{3.1}
\begin{dxtips}
设随机变量 $X$ 取正整数值，则
\[
E[X]=\sum_{k\ge 1}\mathbb P(X\ge k).
\]

\end{dxtips}
\begin{homework}[3.3]
将 $3$ 个球逐个独立且任意地放入编号为 $1$ 至 $4$ 的盒中，
以 $X$ 表示非空盒的最小号码，求 $E(X)$。
\end{homework}
\begin{dxtips}
    这个需要取绝对值再求和
\end{dxtips}
\begin{homework}[3.4]
设离散型随机变量 $X$ 有概率分布
\[
P\!\left(X=(-1)^{i+1}\frac{3^i}{i}\right)=\frac{2}{3^i},\quad i=1,2,\ldots,
\]
说明 $X$ 的期望不存在。
\end{homework}

\begin{homework}[3.7]
设离散型随机变量 $X$ 取正整数的概率依几何数列减少，
选择数列的首项 $a$ 及公比 $q$，使 $E(X)=10$，并计算 $P(X\le 10)$。
\end{homework}

\begin{homework}[3.19]
设二维离散型随机变量 $(X,Y)$ 的概率分布如下：
\[
\begin{array}{c|ccc}
  X\backslash Y & -1 & 0 & 1\\ \hline
  1 & 0.2 & 0.1 & 0.1\\
  2 & 0.1 & 0 & 0.1\\
  3 & 0 & 0.3 & 0.1
\end{array}
\]
求 $E(X)$，$E(Y)$ 与 $E[(X-Y)^2]$。
\end{homework}
\begin{dxtips}
    注意算边缘密度时是对一个确定y算所有不同x的概率所以是12y²从y到1积分。
    但是算E（XY）的时候是先确定x再确定y外层积分就是先确定的。所以x还是取0到1，x取完了y再根据x的取值取值
    \begin{itemize}
  \item \textbf{内层积分：}在外层变量固定时，把另一条方向从起点扫到终点
  （就像秒针在当前这一圈里走完）；
  \item \textbf{外层积分：}把外层变量换一个值，再让内层再扫一遍
  （就像分针动一下，秒针再跑一圈）。
\end{itemize}
\end{dxtips}
\begin{homework}[3.20]
设二维连续型随机变量 $(X,Y)$ 的密度为
\[
f(x,y)=
\begin{cases}
12y^2,& 0\le y\le x\le1,\\[4pt]
0,&\text{其他}.
\end{cases}
\]

求 $E(X)$，$E(Y)$ 和 $E(XY)$。
\end{homework}
\begin{dxtips}
    \[
E(Y)
=\iint y\,f_{X,Y}(x,y)\,dx\,dy
=\int y\left(\int f_{X,Y}(x,y)\,dx\right)dy
=\int y\,f_Y(y)\,dy.
\]
先求y的边缘密度积分对象是x，y相当于常数自然可以乘进去。
我想明白了，其实就是对一个区域积分，这个区域为了满足面积等于 $1$ 每一个的权重都不一样，先算边缘密度比如 $x$ 的边缘密度也就是求出了每一条竖线的值，这个值乘一个数比如 $x$ 和，把这条线拆回无数个点乘 $x$ 是一样的。所以不管是 $E(\,)$ 里面是什么都直接乘就行了。
\end{dxtips}
\begin{note}[二维随机变量的期望及常用情形]
对于二维连续型随机变量 $(X,Y)$，其期望定义推广为
\[
E[g(X,Y)]
=\iint_{\mathbb R^2} g(x,y)\,f(x,y)\,dx\,dy,
\]
其中 $g(x,y)$ 是 $X,Y$ 的某个函数，$f(x,y)$ 为 $(X,Y)$ 的联合概率密度函数。

\vspace{0.5em}
在实际计算中，常用的三种特殊情形为：
\begin{itemize}
  \item 对于 $E(X)$，取 $g(x,y)=x$，于是
  \[
  E(X)=\iint x\,f(x,y)\,dx\,dy.
  \]

  \item 对于 $E(Y)$，取 $g(x,y)=y$，于是
  \[
  E(Y)=\iint y\,f(x,y)\,dx\,dy.
  \]

  \item 对于 $E(XY)$，取 $g(x,y)=xy$，于是
  \[
  E(XY)=\iint xy\,f(x,y)\,dx\,dy.
  \]
\end{itemize}
\end{note}
\begin{dxtips}
    期望对加法和数乘都封闭可以拆开算。
\end{dxtips}
\begin{dxtips}
    若 $X\sim \mathrm{Exp}(\lambda)$，则
\[
E[X]=\frac{1}{\lambda},\qquad \operatorname{Var}(X)=\frac{1}{\lambda^2},
\]
因此
\[
E[X^2]=\operatorname{Var}(X)+\big(E[X]\big)^2=\frac{2}{\lambda^2}.
\]
在这里，$\lambda=1$ 得 $E[X^2]=2$，$\lambda=2$ 得 $E[Y]=\frac{1}{2}$。
\[
E[X]=\int_{0}^{\infty} x\,\lambda e^{-\lambda x}\,dx.
\]
分部积分：取 $u=x\Rightarrow du=dx$，$dv=\lambda e^{-\lambda x}\,dx\Rightarrow v=-e^{-\lambda x}$，则
\[
E[X]=\Big[-x e^{-\lambda x}\Big]_{0}^{\infty}
+\int_{0}^{\infty} e^{-\lambda x}\,dx
=0+\frac{1}{\lambda}
=\frac{1}{\lambda}.
\]
\end{dxtips}
\begin{homework}[3.21]
设随机变量 $X$ 与 $Y$ 相互独立，分别有概率密度函数
\[
f_X(x)=
\begin{cases}
e^{-x}, & x>0,\\
0, & x\le0,
\end{cases}
\qquad
f_Y(y)=
\begin{cases}
2e^{-2y}, & y>0,\\
0, & y\le0.
\end{cases}
\]
求 $E(XY)$ 和 $E(X^2+3Y)$。
\end{homework}
\begin{dxtips}
    正因为二维期望的求法自然期望有线性型应该直接乘x²+3y直接拆成两个积分自然没有问题。
\end{dxtips}
\subsection*{3.2}
\begin{dxtips}
    $S_n=\sum_{k=1}^{n} a_k = a_1\frac{1-q^{\,n}}{1-q}.$
等比数列记忆方法，Sn的和减去Sn的和乘上公比得来的，错位相减只剩第一项1和最后一项q的n次方。容易记忆，开始有n项所以最后一项是第一项乘n-1个公比。
这个和等比级数就是 $q$ 变成 $x$，$a_1=1$，然后因为收敛的所以 $q^n$ 趋近于 $0$ 。对两边关于 $x$ 求导：

\[
\sum_{n=1}^{\infty} n x^{\,n-1} = \frac{1}{(1-x)^2}.
\]

注意：左边的指数是 $n-1$，但我们想要的是 $x^n$ 形式，因此两边同乘以 $x$：

\[
x \sum_{n=1}^{\infty} n x^{\,n-1}
= \sum_{n=1}^{\infty} n x^n
= \frac{x}{(1-x)^2}.
\]

因此得到所需级数：

\[
\boxed{
\sum_{n=1}^{\infty} n x^n = \frac{x}{(1-x)^2},\qquad |x|<1
}
\]
\end{dxtips}
\begin{homework}[3.5]
设离散型随机变量 $X$ 有概率分布 $P(X=n)=2^{-n}$，$n=1,2,\ldots$，
求 $E(X)$ 和 $\mathrm{Var}(X)$。
\end{homework}
\begin{dxtips}
    若随机变量 $X$ 的分布为
\[
P(X=n)=(1-q)q^{n-1},\quad n=1,2,\dots,\quad 0<q<1,
\]
则
\[
E[X]=\frac{1}{1-q},\qquad
\operatorname{Var}(X)=\frac{q}{(1-q)^2}.

\]
\end{dxtips}
\begin{dxtips}
    对于难求的积分一定要变量替换。
\end{dxtips}
\begin{homework}[3.12]
气体分子的速度 $X$ 服从麦克斯韦分布，密度为
\[
f(x)=
\begin{cases}
\dfrac{4}{a^3\sqrt{\pi}}x^2 e^{-x^2/a^2}, & x\ge0,\\
0,& x<0,
\end{cases}
\]
求 $E(X)$ 与 $\mathrm{Var}(X)$。
\end{homework}
\begin{dxtips}
    对于难算的积分观察有没有偶函数奇函数，积分区域是否对称能否简算。
\end{dxtips}
\begin{homework}[3.14]
设随机变量 $X$ 的分布函数为
\[
F(x)=
\begin{cases}
0,& x<-1,\\
0.5+\pi^{-1}\arcsin x,& -1\le x<1,\\
1,& x\ge1,
\end{cases}
\]
求 $E(X)$ 和 $\mathrm{Var}(X)$。
\end{homework}

\begin{homework}[3.16]
设随机变量 $X$ 的期望与方差存在且取值于 $(a,b)$，证明：
\[
a\le E(X)\le b,\qquad 
\mathrm{Var}(X)\le \frac{(b-a)^2}{4}.
\]
\end{homework}
\begin{dxtips}
 \textbf{极坐标变换：}$\quad 
x = r\cos\theta,\quad y = r\sin\theta.$

\textbf{Jacobian 行列式：}

\[
J = \left| \frac{\partial(x,y)}{\partial(r,\theta)} \right|
= \begin{vmatrix}
\cos\theta & -r\sin\theta \\
\sin\theta & r\cos\theta
\end{vmatrix}
= r.
\]

因此：
\[
dx\,dy = r\,dr\,d\theta.
\]

由于 \(x>0,\; y>0\)，仅在第一象限，故：
\[
0 < \theta < \frac{\pi}{2},\qquad r > 0.
\] 
\end{dxtips}
\begin{homework}[3.22]
设二维连续型随机变量 $(X,Y)$ 有概率密度函数
\[
f(x,y)=
\begin{cases}
4xy\exp\{-(x^2+y^2)\}, & x>0,\,y>0,\\[4pt]
0, & \text{其他}.
\end{cases}
\]
求 $Z=\sqrt{X^2+Y^2}$ 的期望和方差。
\end{homework}
\begin{definition}[加马函数]
加马函数定义为
\[
\Gamma(\alpha)=\int_{0}^{+\infty} x^{\alpha-1}e^{-x}\,dx,
\]
并满足递推关系
\[
\Gamma(\alpha+1)=\alpha\,\Gamma(\alpha),\qquad 
\Gamma(n+1)=n!.
\]

其中 $x$ 是积分变量，可以被替换，但必须保证
$x,\ e^{-x},\ dx$ 三者始终对应同一个变量。
被积函数中 $x$ 的幂次为几，则对应的加马函数参数为幂次加 $1$。

\medskip
加马函数的特点是：积分区间为 $[0,+\infty)$，
被积函数前部是 $x$ 的若干次方，后部是 $e^{-x}$。

常用特殊值为
\[
\Gamma(1)=1,\qquad \Gamma\!\left(\tfrac12\right)=\sqrt{\pi}.
\]

\medskip
加马函数还有另一种常用表示形式：
\[
\Gamma(\alpha)=2\int_{0}^{+\infty} x^{2\alpha-1}e^{-x^2}\,dx.
\]

令 $2\alpha-1=n$，即 $\alpha=\frac{n+1}{2}$，则有
\[
\Gamma\!\left(\frac{n+1}{2}\right)
=2\int_{0}^{+\infty} x^{n}e^{-x^2}\,dx.
\]

因此，
\[
\int_{0}^{+\infty} x^{3}e^{-x^2}\,dx
=\frac12\int_{0}^{+\infty} x^{1}e^{-x^2}\,dx
=\frac12\,\Gamma\!\left(\frac{1+1}{2}\right)
=\frac12.
\]
\end{definition}
\begin{homework}[3.23]
设随机变量 $X$ 与 $Y$ 相互独立，$\mathrm{Var}(X)$ 和 $\mathrm{Var}(Y)$ 存在，证明
\[
\mathrm{Var}(XY)
=\mathrm{Var}(X)\mathrm{Var}(Y)
+[E(X)]^2\mathrm{Var}(Y)
+[E(Y)]^2\mathrm{Var}(X).
\]
\end{homework}
\begin{homework}[3.24]
设随机变量 $X\sim B(n,p)$，记 $q=1-p$，证明
\[
E(X-np)^4
=npq(p^3+q^3)+3p^2q^2(n^2-n).
\]
\end{homework}
\subsection*{3.3-3.4}
\begin{homework}[3.25]
设二维离散型随机变量 $(X,Y)$ 有下表所示的概率分布：
\[
\begin{array}{c|ccc}
  X \backslash Y & -1 & 0 & 1 \\ \hline
  -1 & 1/8 & 1/8 & 1/8 \\
   0 & 1/8 & 0   & 1/8 \\
   1 & 1/8 & 1/8 & 1/8
\end{array}
\]
验证 $X$ 与 $Y$ 既不独立，也不相关。
\end{homework}

\begin{homework}[3.26]
设 $A$ 和 $B$ 是有正概率的事件，令
\[
X=
\begin{cases}
1, & A\text{ 出现},\\
0, & A\text{ 不出现},
\end{cases}
\qquad
Y=
\begin{cases}
1, & B\text{ 出现},\\
0, & B\text{ 不出现}.
\end{cases}
\]
若 $X$ 与 $Y$ 不相关，证明 $X$ 与 $Y$ 相互独立。
\end{homework}
\begin{definition}[相关系数与协方差矩阵]
相关系数的定义为
\[
\rho_{X,Y}
=\frac{\operatorname{Cov}(X,Y)}
{\sqrt{\operatorname{Var}(X)\operatorname{Var}(Y)}}.
\]

这是相关系数最根本的形式，适用于任意具有有限方差的随机变量。
另一方面，二维随机向量 $(X,Y)$ 的协方差矩阵为
\[
\Sigma=
\begin{pmatrix}
E\!\big[(X-E[X])^2\big] & E\!\big[(X-E[X])(Y-E[Y])\big] \\
E\!\big[(X-E[X])(Y-E[Y])\big] & E\!\big[(Y-E[Y])^2\big]
\end{pmatrix}.
\]

其中对角线元素给出方差，非对角线元素给出协方差。
将相关系数写成协方差矩阵的形式，可得
\[
\rho_{X,Y}
=\frac{\Sigma_{12}}{\sqrt{\Sigma_{11}\Sigma_{22}}}.
\]

因此，相关系数实际上就是协方差在对应方差尺度下的标准化结果。
\end{definition}
\begin{homework}[3.28]
在 $10$ 件产品中有 $2$ 件一级品，$7$ 件二级品，$1$ 件次品。从中不放回地抽取 $3$ 件，
用 $X$ 表示抽到的一级品数，$Y$ 表示抽到的二级品数，求 $X$ 与 $Y$ 的相关系数。
\end{homework}

\begin{homework}[3.30]
设二维连续型随机变量 $(X,Y)$ 有概率密度函数
\[
f(x,y)=
\begin{cases}
1, & |y|<x,\ 0<x<1,\\[4pt]
0, & \text{其他}.
\end{cases}
\]
求 $\mathrm{Cov}(X,Y)$。
\end{homework}

\begin{homework}[3.32]
设 $X,Y,Z$ 是随机变量，且
\[
E(X)=E(Y)=1,\quad E(Z)=-1,\quad 
\mathrm{Var}(X)=\mathrm{Var}(Y)=\mathrm{Var}(Z)=1,
\]
\[
\rho_{X,Y}=0,\quad \rho_{X,Z}=0.5,\quad \rho_{Y,Z}=-0.5.
\]
求 $E(X+Y+Z)$ 和 $\mathrm{Var}(X+Y+Z)$。
\end{homework}

\begin{homework}[3.34]
设 $n>m$，随机变量 $X_1,X_2,\ldots,X_{n+m}$ 相互独立，且具有相同的分布函数，
$\mathrm{Var}(X_1)$ 存在。在
\[
Y=\sum_{i=1}^{n} X_i,\qquad Z=\sum_{k=1}^{m} X_{n+k}
\]
下，求 $\rho_{Y,Z}$。
\end{homework}
\subection*{3.5}
\begin{homework}[3.37]
设每天进入货站的货物件数 $N$ 有如下概率分布：
\[
\begin{array}{c|cccccc}
N=n & 10 & 11 & 12 & 13 & 14 & 15 \\ \hline
P(N=n) & 0.05 & 0.10 & 0.10 & 0.20 & 0.35 & 0.20
\end{array}
\]
并设每天到达的货物中次品的概率均为 $0.10$。
以 $X$ 表示每天进货中次品的件数，求 $E(X)$。
\end{homework}

\begin{homework}[3.38]
设随机变量 $X_1 \sim P(\lambda_1)$，$X_2 \sim P(\lambda_2)$，且二者独立。

\begin{enumerate}[(1)]
  \item 证明：当 $n \ge 2$ 时，有
  \[
  P(X_1=i \mid X_1+X_2=n)
  = C_n^i
    \left(\frac{\lambda_1}{\lambda_1+\lambda_2}\right)^i
    \left(\frac{\lambda_2}{\lambda_1+\lambda_2}\right)^{n-i},
  \qquad i=0,1,2,\ldots,n.
  \]
  \item 求 $E(X_1\mid X_1+X_2=n)$ 和 $E(X_2\mid X_1+X_2=n)$。
\end{enumerate}
\end{homework}

\begin{homework}[3.39]
设二维连续型随机变量 $(X,Y)$ 有概率密度函数
\[
f(x,y)=
\begin{cases}
8xy, & 0<y<x<1,\\[4pt]
0, & \text{其他}.
\end{cases}
\]
求条件期望 $E(X\mid Y=y)$ 与 $E(Y\mid X=x)$。
\end{homework}
\section{第四章}
\subsection*{4.1 }
\begin{homework}[4.1]
\begin{dxtips}[]
只要是求特征函数的导数，都可以这样理解：因为 $t=0$ 的存在，
所有含有 $e^{it}$ 的项在 $t=0$ 处最终都等于 $1$。
可以把 $e^{it}$ 以及它前面的系数整体看作关于 $t$ 的一个函数。

从本质上来说，求完特征函数以后，$e^{itx}$ 中的 $x$ 一定会消失。
因为本来就是要把 $x$ 代入具体的数值然后全部加和，
这个变量在期望运算完成后不可能再显式出现，
而是被编码进特征函数这一关于 $t$ 的函数中。

然后由于 $e^{it}$ 不管它前面的系数是多少，
对 $t$ 求一次导数都会下来一个 $i$，
因此特征函数的导数非常好算，
这也正是特征函数在计算期望和方差时
能够大大简化计算的原因。
    令 $g(t)=1-qe^{it}$，则 $\varphi(t)=p\,g(t)^{-1}$，且
\[
g'(t)=-iqe^{it},\qquad g''(t)=qe^{it}.
\]
由 $(g^{-1})'=-g'/g^2$ 与 $(g^{-1})''=\dfrac{2(g')^2-g''g}{g^3}$，得
\[
\varphi'(t)=p\,(g^{-1})'=-p\,\frac{g'(t)}{g(t)^2},
\qquad
\varphi''(t)=p\,(g^{-1})''
=p\,\frac{2(g'(t))^2-g''(t)g(t)}{g(t)^3}.
\]
代入 $t=0$（此时 $g(0)=1-q=p$）可得
\[
\varphi'(0)=\frac{iq}{p},
\qquad
\varphi''(0)=-\,\frac{q(1+q)}{p^2}.
\]
\end{dxtips}
设随机变量 $X$ 服从几何分布，有概率分布
\[
P(X=k)=pq^{k},\qquad k=0,1,2,\ldots,
\]
其中 $0<p<1,\; q=1-p$. 求 $X$ 的特征函数、$E(X)$ 和 $\operatorname{Var}(X)$。
\end{homework}

\begin{homework}[4.2]
设随机变量 $X$ 服从帕斯卡分布，有概率分布
\[
P(X=k)=C_{k-1}^{\,r-1} p^{\,r} q^{\,k-r},\qquad k=r,r+1,r+2,\ldots,
\]
其中 $0<p<1,\; q=1-p,\; r$ 为正整数。求 $X$ 的特征函数。
\end{homework}
\subsection*{4.2}
\begin{dxtips}
    先弄出有e^itk的项先\[
e^{i\theta}=\cos\theta+i\sin\theta.
\]
欧拉公式这个模长是1除了i就是要放进sin，cos里面的
。cos是落在x轴上的所以是实数，sin落在y轴上所以是虚数
\end{dxtips}
\begin{dxtips}
   看见 $e^{-ikt}$ 也不要慌，它前面的系数只是随机变量取值为 $x=-k$ 时的概率罢了。
把符号不同的 $e^{ikt}$ 分开来看：正号对应 $x=k$，负号对应 $x=-k$。
只要这些系数都非负，并且最终加和为 $1$，就能把该函数写成
\[
\varphi(t)=\sum_{x\in\mathbb Z} P(X=x)\,e^{itx}=E\!\big(e^{itX}\big),
\]
从而证明它是某个离散型随机变量的特征函数。
\end{dxtips}
\begin{homework}[4.5]
设 $a_k\ge 0,\ \sum_{k=0}^\infty a_k=1$，证明
\[
\varphi_1(t)=\sum_{k=0}^\infty a_k \cos(kt), 
\qquad 
\varphi_2(t)=\sum_{k=0}^\infty a_k e^{jkt}
\]
均为特征函数，并求相应的离散型随机变量的概率分布（傅里叶逆变换）。
\end{homework}
\begin{dxtips}
    经历过福利呀变换自然有eit把这项弄出来是关键。
\end{dxtips}
\begin{homework}[4.6]
设特征函数分别为 $\cos t$ 和 $\cos^2 t$，求相应的离散型随机变量的概率分布（傅里叶逆变换）。
\end{homework}
\begin{dxtips}
连续切单调他就服从0到1均匀分布
\textbf{(1) } $Y=aF(X)+b$

令 $U=F(X)$，则 $Y=aU+b$。按特征函数的定义，
\[
\varphi_Y(t)=E\!\left(e^{it(aU+b)}\right)
=e^{itb}\,E\!\left(e^{iatU}\right)
=e^{itb}\int_{0}^{1}e^{iatu}\,du.
\]
因此
\[
\varphi_{aF(X)+b}(t)=
\begin{cases}
e^{itb}\,\dfrac{e^{iat}-1}{iat}, & t\neq 0,\\[1ex]
1, & t=0.
\end{cases}
\]
\end{dxtips}
\begin{homework}[4.7]
设 $a,b$ 是实常数，$a\neq 0$，随机变量 $X$ 的分布函数 $F(x)$ 连续且严格单调递增，分别求 $aF(X)+b$ 及 $\ln F(X)$ 的特征函数。
\end{homework}
\subsection*{4.3-4.4}
\begin{homework}[4.11]
设随机变量 $X_1,X_2,\ldots,X_n$ 独立同分布，且均服从标准正态分布，证明  
\[
n^{-1/2}\sum_{i=1}^n X_i
\]
服从标准正态分布。
\end{homework}

\begin{homework}[4.12]
设 $(X_1,X_2)$ 服从非退化的二维正态分布  
\[
N\!\left(0,\sigma_1^2; \; 0,\sigma_2^2;\; \frac{\sigma_{12}}{\sigma_1\sigma_2}\right),
\]
其中 $\sigma_{12}=\operatorname{Cov}(X_1,X_2)$。求  
\[
E\bigl[(X_1^2-\sigma_1^2)(X_2^2-\sigma_2^2)\bigr].
\]
\end{homework}
\subsection{4.5}
\begin{homework}[4.14]
设随机变量 $X$ 的母函数如下，求 $X$ 的概率分布：  
(1)\ $0.25(1+s)^2$；  
(2)\ $\displaystyle \frac{1}{2-s}$；  
(3)\ $e^{\,s-1}$。
\end{homework}

\begin{homework}*[4.15]
在伯努利试验中，假设每次试验成功的概率为 $p$，如果记 $X$ 为第 $r$ 次失败之前成功出现的总次数，求 $X$ 的概率分布、母函数、期望和方差。
\end{homework}

\begin{homework}[4.16]
设 $\{X_k,\,k\ge 1\}$ 是独立随机变量序列，且
\[
P(X_k=i)=\frac{1}{m},\qquad i=0,1,2,\ldots,m-1,\quad k=1,2,\ldots,
\]
令 $S_n=\sum_{k=1}^{n} X_k$，求 $S_n$ 的母函数。
\end{homework}
\chapter{第五章}
\section{5.1}
\begin{dxtips}
    求期望不用管相关性，只用知道骰子点数和骰子概率就行了。但是求方差就得注意相关性
    \[
\operatorname{Var}\!\left(\sum_i X_i\right)
=\sum_i \operatorname{Var}(X_i)
+2\sum_{i<j}\operatorname{Cov}(X_i,X_j).
\]
\end{dxtips}
\textbf{（1）两个变量的例子}

设 $a=X_1-\mu_1$，$b=X_2-\mu_2$，则
\[
(a+b)^2=a^2+b^2+2ab.
\]
对应回去有
\[
\big((X_1-\mu_1)+(X_2-\mu_2)\big)^2
=(X_1-\mu_1)^2+(X_2-\mu_2)^2
+2(X_1-\mu_1)(X_2-\mu_2).
\]
其中系数 $2$ 来自交叉项 $ab$ 出现了两次：一次来自 $a\cdot b$，一次来自 $b\cdot a$。
\begin{homework}[5.3]
将编号为 $1,2,\ldots,n$ 的球放入编号为 $1,2,\ldots,n$ 的盒，每盒只放一球，
以 $S_n$ 表示球与盒的编号正好相同的个数，证明
\[
\frac1n \bigl[S_n - E(S_n)\bigr] \xrightarrow{P} 0 
\qquad (n\to\infty).
\]
\end{homework}

%------------------------------------------------------
\begin{theorem}[Kolmogorov 判别法]
设 $\{X_n\}$ 是独立随机变量序列，各 $\operatorname{Var}(X_n)=\sigma_n^2$ 存在，且
\[
\sum_{n=1}^{\infty}\frac{\sigma_n^2}{n^2}<\infty,
\]
则随机变量序列 $\{X_n\}$ 服从强大数定律。
和切比雪夫不一样切比雪夫下面是确定ε²和n²，n趋紧于无穷，但是这个下面是n²的求和
\end{theorem}

\begin{homework}[5.4]
设 $\{X_n\}$ 为独立随机变量序列，满足
\[
P(X_n = n^\lambda) = P(X_n = -n^\lambda) = 0.5 , \qquad n=1,2,\ldots,
\]
其中常数 $\lambda<0.5$。证明 $\{X_n\}$ 服从大数定律。
\end{homework}

%------------------------------------------------------

\begin{homework}[5.5]
设 $\{X_n\}$ 为独立同分布随机变量序列，$X_1$ 的分布函数为
\[
F(x)=\frac12 + \frac1\pi \arctan x , \qquad -\infty < x < \infty .
\]
问该随机变量序列是否服从大数定律？并说明结论。
\end{homework}

%------------------------------------------------------
\begin{dxtips}
    一般先求期望期望=0，然后用切比雪夫不等式连接X-E（X）与方差σ²。
\end{dxtips}
\begin{homework}[5.7]
设独立随机变量序列 $\{X_n\}$ 分别有以下概率分布，讨论其是否服从大数定律：

\begin{itemize}
  \item[(1)] $P(X_n = 2^n)=P(X_n = -2^n)=0.5$；
  \item[(2)] $P(X_n = 2^n)=P(X_n = -2^n)=2^{-(2n+1)}$, 
              \quad $P(X_n=0)=1-2^{-2n}$；
  \item[(3)] $P(X_n = n)=P(X_n = -n)=0.5 n^{-0.5}$, 
              \quad $P(X_n=0)=1-n^{-0.5}$。
\end{itemize}
\end{homework}

%------------------------------------------------------

\begin{homework}[5.8]
设 $\{X_n\}$ 为随机变量序列，且 $i\ne j$ 时协方差满足 $\operatorname{Cov}(X_i,X_j)<0$。
讨论该随机变量序列是否服从大数定律。
\end{homework}

%------------------------------------------------------

\begin{homework}[5.9]
设 $\{X_n\}$ 是随机变量序列，各 $\operatorname{Var}(X_n)$ 存在，且 $\operatorname{Var}(X_n)\le c$，
其中常数 $c$ 与 $n$ 无关。又设当 $|i-j|\to\infty$ 时，相关系数 $\rho_{ij}\to 0$。
讨论 $\{X_n\}$ 是否服从大数定律。
\end{homework}
\section{5.2}
中心极限定理就是说：只要你知道独立同分布（或满足条件的）随机变量的均值 $E[X]$ 和方差 $\operatorname{Var}(X)$，就能把大量样本的和或均值近似看成正态分布，从而用正态分布表计算关于“和/平均数”的概率，或反求所需的样本量。
\begin{dxtips}
    先算x的期望E（x）和x的方差var（x） 然后得到Sn的期望nE（x），和Sn的方差nvar（x），此时他服从的是（nμ，nvar（x）分布。要求他小于一个数字的概率，左边右边都归一化处理我们要求的是P（X小于a）所以对左边操作也得对右边操作。归一化了以后就可以查表了。
\end{dxtips}
\begin{homework}[5.14]
用计算机做统计模拟时，要求数据准确到小数点后第 $6$ 位，而对第 $6$ 位后面的数字作“四舍五入”处理。设所有舍入误差相互独立，且均服从
\[
U(-0.5\times 10^{-6},\,0.5\times 10^{-6}).
\]
\begin{itemize}
  \item[(1)] 将 $1500$ 个数相加，求误差总和的绝对值超过 $15\times 10^{-6}$ 的概率；
  \item[(2)] 求最多多少个数相加，可使误差总和的绝对值小于 $10^{-6}$ 的概率不小于 $0.9$。
\end{itemize}
\end{homework}

\begin{solution}
设舍入误差为 $X_i\sim U(-a,a)$，其中 $a=0.5\times 10^{-6}$，并令
$S_n=\sum_{i=1}^n X_i$。
则
\[
E[X_i]=0,\qquad 
\operatorname{Var}(X_i)=\frac{a^2}{3}
=\frac{(0.5\times 10^{-6})^2}{3}
=8.3333\times 10^{-14}.
\]
因此
\[
E[S_n]=0,\qquad 
\operatorname{Var}(S_n)=n\operatorname{Var}(X_i)=\frac{na^2}{3},\qquad
\sigma_n=\sqrt{\operatorname{Var}(S_n)}
=\sqrt{\frac{na^2}{3}}.
\]
由中心极限定理，$S_n\approx N(0,\sigma_n^2)$。

\begin{itemize}
  \item[(1)] 当 $n=1500$ 时，
  \[
  \sigma_{1500}
  =\sqrt{1500\times 8.3333\times 10^{-14}}
  =\sqrt{1.25\times 10^{-10}}
  =1.1180\times 10^{-5}.
  \]
  因而
  \[
  P\big(|S_{1500}|>15\times 10^{-6}\big)
  \approx
  2\Bigl(1-\Phi\Bigl(\frac{15\times 10^{-6}}{1.1180\times 10^{-5}}\Bigr)\Bigr)
  =2\bigl(1-\Phi(1.342)\bigr)\approx 0.180.
  \]

  \item[(2)] 要求
  \[
  P(|S_n|<10^{-6})
  \approx
  2\Phi\Bigl(\frac{10^{-6}}{\sigma_n}\Bigr)-1
  \ge 0.9
  \iff
  \frac{10^{-6}}{\sigma_n}\ge z_{0.95}=1.645.
  \]
  又
  \[
  \sigma_n=\frac{\sqrt{n}\,a}{\sqrt{3}},
  \]
  故
  \[
  \frac{10^{-6}}{\sigma_n}
  =
  \frac{10^{-6}}{\sqrt{n}\,a/\sqrt{3}}
  =
  \frac{\sqrt{3}}{\sqrt{n}}\cdot
  \frac{10^{-6}}{0.5\times 10^{-6}}
  =
  \frac{3.464}{\sqrt{n}}.
  \]
  因而
  \[
  \frac{3.464}{\sqrt{n}}\ge 1.645
  \iff
  \sqrt{n}\le 2.105
  \iff
  n\le 4.43,
  \]
  所以最多取 $n_{\max}=4$。
\end{itemize}
\end{solution}

\begin{theorem}[De Moivre--Laplace 中心极限定理]
设随机变量 $X\sim B(n,p)$，记 $q=1-p$，则对任意实数 $x$，有
\[
\lim_{n\to+\infty}
P\!\left\{
\frac{X-np}{\sqrt{npq}}\le x
\right\}
=
\frac{1}{\sqrt{2\pi}}
\int_{-\infty}^{x} e^{-t^{2}/2}\,dt .
\]
\end{theorem}
\begin{homework}[5.15]
某单位的电话总机共有 $200$ 部分机，设各分机是否使用外线相互独立，各时刻每部分机使用外线的概率均为 $0.05$。求总机至少需要多少条外线，才能以 $90\%$ 以上的概率使各分机在使用外线时不必等待？
\end{homework}

\begin{solution}
令 $S$ 为同一时刻使用外线的分机数，则
\[
S\sim \Bin(200,0.05),\qquad \mu=E[S]=200\times 0.05=10,\qquad
\sigma^2=\Var(S)=200\times 0.05\times 0.95=9.5,
\]
\[
\sigma=\sqrt{9.5}=3.082.
\]
若外线数为 $m$，则“不等待”等价于 $S\le m$。用正态近似并作连续性修正：
\[
P(S\le m)\approx \Phi\Big(\frac{m+0.5-\mu}{\sigma}\Big)\ge 0.9
\iff \frac{m+0.5-10}{3.082}\ge z_{0.9}=1.282.
\]
因此
\[
m\ge 10-0.5+3.082\times 1.282=13.45,
\]
取最小整数得 $m_{\min}=14$。
\end{solution}
\begin{dxtips}
    这个 $0.5$ 来自连续性修正（continuity correction）。

因为 $S\sim \text{Bin}(200,0.05)$ 是离散整数随机变量（只能取 $0,1,2,\dots$），而正态分布是连续分布。为了使正态近似更加准确，通常将
\[
P(S\le m)
\]
改写为连续型随机变量的概率。

具体地，事件 $S\le m$ 等价于 $S<m+1$，对整数型随机变量而言，其对应的连续近似区间为
\[
S\le m \quad \approx \quad Y\le m+0.5,
\]
其中 $Y$ 表示用正态分布近似 $S$ 的连续随机变量。

直观地看，整数点 $m$ 在连续数轴上对应区间 $[m-0.5,\;m+0.5)$，因此事件“$S\le m$”的连续边界应取为 $m+0.5$。

于是有正态近似公式
\[
P(S\le m)\approx 
\Phi\!\left(\frac{m+0.5-\mu}{\sigma}\right).
\]
\end{dxtips}

\begin{homework}[5.16]
在一批用于建筑的木材中，$80\%$ 的长度不小于 $4\,\mathrm m$。从中任取 $100$ 根，求这 $100$ 根中至少有 $30$ 根短于 $4\,\mathrm m$ 的概率。
\end{homework}

\begin{solution}
令 $X$ 为短于 $4\,\mathrm m$ 的根数，则
\[
P(\text{短于 }4\,\mathrm m)=0.2,\qquad X\sim \Bin(100,0.2),
\]
\[
\mu=E[X]=100\times 0.2=20,\qquad \sigma^2=\Var(X)=100\times 0.2\times 0.8=16,\qquad \sigma=4.
\]
所求为 $P(X\ge 30)$。用正态近似并作连续性修正：
\[
P(X\ge 30)=P(X\ge 29.5)\approx 1-\Phi\Big(\frac{29.5-20}{4}\Big)
=1-\Phi(2.375)\approx 0.0088.
\]
\end{solution}
